\section{Fundamentação teórica}

A segurança ofensiva de redes Wi-Fi compreende um conjunto de técnicas e metodologias destinadas a identificar, explorar e relatar vulnerabilidades em sistemas e infraestruturas de rede, com o propósito proativo de fortalecer suas defesas. Diferente da segurança defensiva, que visa proteger sistemas contra ataques, a abordagem ofensiva simula as ações de adversários reais para avaliar a resiliência dos mecanismos de proteção. A área da segurança ofensiva vem ganhando atenção por se mostrar eficiente e ajudar no treinamento de equipes e refinamento nas técnicas de detecção e prevenção de ataques.

Neste projeto, a segurança ofensiva é aplicada por meio de uma ferramenta que realiza varreduras, enumeração e ataques em ambientes sem fio. Técnicas como flooding (inundação) de beacons, deautenticação de dispositivos pelo protocolo IEEE 802.11, spoofing (falsificação) de endereços IP/MAC são fundamentais para o teste de penetração, explorando falhas em protocolos, configurações inadequadas ou limitações de implementação.

Para a implementação de uma ferramenta ofensiva eficaz, é essencial o domínio dos protocolos de rede que atuam em diversas camadas do modelo TCP/IP, em especial na camada 2, onde há a presença do protocolo IEEE 802.11. O ICMP (\textit{Internet Control Message Protocol}), utilizado para diagnóstico, serve de base para enumeração e busca de informações de dispositivos. O padrão IEEE 802.11 (Wi-Fi) é essencial, sendo explorado por meio da captura de beacons e ataques de desautenticação, que manipulam quadros de gerenciamento.